% nff-power-meter — experiment datasheet
%
% Build:  pdflatex datasheet.tex   (run twice for the ToC / references)
% Deps:   circuitikz, booktabs, siunitx, xcolor, geometry, fancyhdr, enumitem
%         All are in a standard TeX Live / MiKTeX "full" install.

\documentclass[11pt,a4paper]{article}

\usepackage[margin=22mm]{geometry}
\usepackage{booktabs}
\usepackage{siunitx}
\usepackage{circuitikz}
\usepackage{xcolor}
\usepackage{fancyhdr}
\usepackage{enumitem}
\usepackage{amsmath}
\usepackage{tabularx}
\usepackage[hidelinks]{hyperref}

\sisetup{
  per-mode              = symbol,
  range-units           = single,
  detect-weight         = true,
  group-minimum-digits  = 4,
}

\definecolor{warn}{RGB}{176,0,32}
\definecolor{ok}{RGB}{0,120,60}
\providecommand{\good}[1]{\textcolor{ok}{\textbf{#1}}}
\providecommand{\bad}[1]{\textcolor{warn}{\textbf{#1}}}
\newcommand{\yes}{\textcolor{ok}{\textbf{yes}}}
\newcommand{\no}{\textcolor{warn}{\textbf{no}}}

\pagestyle{fancy}
\fancyhf{}
\lhead{\small nff-power-meter}
\chead{\small Experiment datasheet}
\rhead{\small Rev.\ A --- 2026-07-13}
\cfoot{\small \thepage}

\title{\textbf{nff-power-meter}\\[2mm]
\large Shunt-based energy measurement of ESP32 firmware operations\\[1mm]
\normalsize Experiment datasheet --- Revision A}
\author{Bench rig: STM32 Nucleo-F446RE + ESP32 DevKit}
\date{2026-07-13}

\begin{document}
\maketitle
\thispagestyle{fancy}

\begin{abstract}
\noindent
This document specifies the bench rig used to measure the electrical energy consumed
by an ESP32 device while it executes an \texttt{nff} operation --- principally
\texttt{nff ota}, the only \texttt{nff} command that places a sustained load
(WiFi RX + flash writes) on the target. The measurand is \textbf{joules per OTA}.
A \SI{1}{\ohm} shunt in the ESP32's ground return is sampled at \SI{200}{kSps} by an
STM32 Nucleo-F446RE, which accumulates charge and energy \emph{in firmware} and reports running
totals on demand. The host never integrates, and therefore cannot silently drop samples and
under-report energy.

\medskip
\noindent
Three guards decide whether a run is reportable at all, and each exists because the rig produced
a confident, plausible, \emph{wrong} number during bring-up (\S\ref{sec:honesty}): a DMA-overrun
counter (dropped conversions), a window cross-check (a lost \texttt{ZERO} once turned a
\SI{30}{\second} run into \SI{91}{\second} and \SI{26}{\joule}), and an active continuity probe
(a rig wired to \emph{nothing} reported \SI{590}{\milli\ampere} and a \SI{4.94}{\volt} rail). If
any fails, no joules figure is reported.
\end{abstract}

\clearpage
\tableofcontents
\clearpage

% ---------------------------------------------------------------------------
\section{Scope and measurand}

\begin{description}[style=nextline,leftmargin=0pt]
  \item[Primary measurand] Marginal energy $E_\text{marg}$ (\si{\joule}) consumed by the
        ESP32 module over and above its idle draw, for the duration of one \texttt{nff}
        command.
  \item[Secondary] Mean current, peak current, supply voltage, total charge, run duration.
  \item[Purpose] (a) reason about battery life; (b) provide a regression gate
        (\texttt{-{}-max-joules}) so an agent can detect a firmware change that increases
        power draw.
\end{description}

\noindent
The device under test (DUT) is the \emph{whole ESP32 devkit board}, not the ESP32 die.
The measurement therefore includes the onboard LDO loss, the USB--UART bridge, and the
power LED. These are constant, so they cancel in the baseline-subtracted marginal figure
(\S\ref{sec:math}).

% ---------------------------------------------------------------------------
\section{Measurement principle}

The shunt is placed in the DUT's \textbf{ground return}, not its supply line. The voltage
developed across it is therefore \emph{ground-referenced} and can be read single-ended by a
plain ADC input. This is the sole reason the rig needs no instrumentation amplifier.

Because the shunt is low-side, the DUT's local ground (node \texttt{E\_GND}) floats above
system ground by $I \cdot R_s$ --- up to \SI{300}{\milli\volt} at \SI{300}{\milli\ampere}.
The voltage actually seen \emph{across the DUT's terminals} is consequently the rail voltage
minus that lift, and it is that difference --- not the rail alone --- that must multiply the
current to obtain true power.

% ---------------------------------------------------------------------------
\section{Schematic}

\subsection{Power path}

\begin{center}
\begin{circuitikz}[american, scale=0.92, transform shape]

  % ===== Nucleo block ====================================================
  \draw[thick] (0,-6.2) rectangle (3,1.1);
  \node[align=center] at (1.5,-2.5) {\small Nucleo\\[-0.5mm]\small F446RE};
  \node[align=center,anchor=north] at (1.5,-4.6)
        {\scriptsize USB $\rightarrow$ PC\\[-0.6mm]\scriptsize (only cable)};

  % Nucleo pin stubs
  \draw (3,0.5)  to[short,-*] (3.4,0.5)  ; \node[anchor=south west] at (3.45,0.55) {\scriptsize\texttt{5V}};
  \draw (3,-1)   to[short,*-] (3.4,-1)   ; \node[anchor=south west] at (3.05,-0.92){\scriptsize\texttt{PA1}/A1};
  \draw (3,-4.2) to[short,*-] (3.4,-4.2) ; \node[anchor=south west] at (3.05,-4.12){\scriptsize\texttt{PA0}/A0};
  \draw (3,-5.5) to[short,*-] (3.4,-5.5) ; \node[anchor=south west] at (3.05,-5.42){\scriptsize\texttt{GND}};

  % ===== supply rail =====================================================
  \draw (3.4,0.5) to[short] (12,0.5);

  % ===== rail divider (10k/10k) =========================================
  \draw (5.5,0.5) to[R={$10\,\mathrm{k}$}, *-] (5.5,-1);
  \draw (5.5,-1)  to[short,*-] (3.4,-1);                    % tap -> PA1
  \draw (5.5,-1)  to[R={$10\,\mathrm{k}$}] (5.5,-2.6);
  \draw (5.5,-2.6) to[short] (5.5,-3.1) node[ground]{};

  % optional S/H buffer cap on PA1 (label flipped to the left)
  \draw (4.2,-1) to[C, l_={$100\,\mathrm{nF}$}, *-] (4.2,-2.4) node[ground]{};

  % ===== optional bulk cap: VIN -> E_GND (NOT to N_GND) ==================
  \draw (10,0.5) to[C, l_={$100\,\mu\mathrm{F}$}, *-*] (10,-4.2);
  \node[anchor=east] at (9.75,-3.3) {\scriptsize optional};

  % ===== DUT =============================================================
  \draw[thick] (11,-3) rectangle (15,0);
  \node[align=center] at (13,-1.5) {\small ESP32 devkit\\[-0.5mm]\scriptsize (DUT)};
  \draw (12,0)  to[short,*-] (12,0.5);                       % VIN up to rail
  \node[anchor=south west] at (12.05,0.05) {\scriptsize\texttt{VIN}};
  \draw (12,-3) to[short] (12,-4.2);                         % GND down to E_GND
  \node[anchor=north west] at (12.05,-3.05) {\scriptsize\texttt{GND}};
  \node[anchor=west,align=left] at (15.15,-1.5) {\scriptsize\textbf{USB:}\\[-0.6mm]\scriptsize not connected};

  % ===== E_GND rail (DUT-local ground, floats above N_GND) ===============
  \draw (8,-4.2) to[short,*-*] (12,-4.2);
  \node[anchor=south] at (8.7,-4.08) {\scriptsize\texttt{E\_GND}};

  % ===== shunt sense to PA0 ==============================================
  \draw (8,-4.2) to[short] (3.4,-4.2);

  % ===== shunt ===========================================================
  \draw (8,-4.2) to[R={$R_s=1\,\Omega$}] (8,-5.5);

  % ===== N_GND rail ======================================================
  \draw (8,-5.5) to[short,*-] (3.4,-5.5);
  \draw (6,-5.5) to[short,*-] (6,-6.1) node[ground]{};
  \node[anchor=north] at (7.1,-5.70) {\scriptsize\texttt{N\_GND} (ADC reference)};

\end{circuitikz}
\end{center}

\noindent
\textbf{Read it as a loop.} Current leaves the Nucleo \texttt{5V} pin, enters the DUT at
\texttt{VIN}, passes through the entire board, exits at the DUT \texttt{GND} pin, and returns
to \texttt{N\_GND} through $R_s$. Every milliampere the DUT draws is forced through the shunt.
\emph{The DUT must have no second ground path} --- see \S\ref{sec:rules}.

\medskip
\noindent
\fcolorbox{warn}{white}{\parbox{0.97\textwidth}{%
\textbf{The return-path trap.} The bulk capacitor $C_b$ returns to \texttt{E\_GND} --- the DUT's
\emph{own} ground, at the top of the shunt --- \textbf{not} to \texttt{N\_GND}. A component
returned to \texttt{N\_GND} carries DUT current that bypasses the shunt entirely and is therefore
invisible to the meter: a silent, systematic under-read. The same rule governs the marker
pulldown below. Note $C_f$ (on \texttt{PA1}) is exempt: it carries no DUT current, only ADC
sample-and-hold charge, so it returns to \texttt{N\_GND}.}}

\subsection{Phase marker (Phase 5, optional)}

Drawn separately because its ground return is the easiest thing on this rig to get wrong.

\begin{center}
\begin{circuitikz}[american, scale=0.92, transform shape]

  % DUT
  \draw[thick] (0,-1.2) rectangle (3,0.6);
  \node at (1.5,-0.3) {\small ESP32 (DUT)};
  \node[anchor=south west] at (3.05,-0.35) {\scriptsize\texttt{GPIO27}};
  \node[anchor=north east] at (1.45,-1.25) {\scriptsize\texttt{GND}};

  % GPIO -> 1k -> PA10
  \draw (3,-0.5) to[R, l_={$1\,\mathrm{k}$}, *-*] (5.6,-0.5);
  \draw (5.6,-0.5) to[short,-o] (7.6,-0.5) node[right]{\small\texttt{PA10} / D2 (5\,V-tolerant)};

  % pulldown -> E_GND
  \draw (5.6,-0.5) to[R={$10\,\mathrm{k}$}] (5.6,-2.8);
  \draw (5.6,-2.8) to[short] (1.2,-2.8);

  % DUT GND down to E_GND node
  \draw (1.2,-1.2) to[short,-*] (1.2,-2.8);
  \node[anchor=south east] at (1.15,-2.75) {\scriptsize\texttt{E\_GND}};

  % shunt down to N_GND
  \draw (1.2,-2.8) to[R={$R_s$}] (1.2,-4.3);
  \draw (1.2,-4.3) to[short] (1.2,-4.7) node[ground]{};
  \node[anchor=west] at (1.35,-4.55) {\scriptsize\texttt{N\_GND}};

  % annotation
  \draw[->,thick,warn] (5.0,-3.1) .. controls (4.0,-3.9) and (3.0,-3.5) .. (2.0,-3.2);
  \node[anchor=west,warn,align=left] at (5.1,-3.35)
        {\scriptsize\textbf{pulldown returns to \texttt{E\_GND},}\\[-0.6mm]
         \scriptsize\textbf{above the shunt --- never \texttt{N\_GND}}};

\end{circuitikz}
\end{center}

\noindent
Every ESP32 GPIO goes high-impedance during reset, and an OTA \emph{ends} in a reboot --- so
without the pulldown the marker line floats through precisely the phase transition being
captured. The \SI{10}{\kilo\ohm} holds it at a defined low; the \SI{1}{\kilo\ohm} protects both
pins should one end be misconfigured as an output. Returning the pulldown to \texttt{N\_GND}
would route $\approx\SI{330}{\micro\ampere}$ of DUT current around the shunt.

\medskip
\noindent
\textbf{Two markers are better than one.} A single pin distinguishes only ``in a phase / not in a
phase''. Two pins (e.g.\ GPIO27 + GPIO26 into two FT inputs) give a 2-bit code --- exactly the four
OTA phases worth separating: download / write / verify / reboot. Same wiring cost.

% ---------------------------------------------------------------------------
\section{Bill of materials}

\begin{center}
\begin{tabularx}{\textwidth}{llX}
\toprule
\textbf{Ref} & \textbf{Value} & \textbf{Notes} \\
\midrule
MCU     & Nucleo-F446RE            & Meter + host link. Only device connected to the PC. \\
DUT     & ESP32 devkit (WROOM-32)  & Powered exclusively from the Nucleo \SI{5}{\volt} pin. \\
$R_s$   & \SI{1}{\ohm}, \SI{0.25}{\watt} & Shunt. Ten \SI{10}{\ohm} in parallel is an acceptable substitute (also averages tolerance down). \\
$R_1,R_2$ & \SIrange{1}{10}{\kilo\ohm} & Rail divider, matched pair; \SI{1}{\percent} preferred (ratio error is a direct gain error on energy). \SI{10}{\kilo\ohm}/\SI{10}{\kilo\ohm} settles with $\approx6\times$ margin --- see \S\ref{sec:adc}. Ratio is \num{2.0} either way. \\
$C_b$   & \SI{100}{\micro\farad} & \emph{Optional.} Bulk decoupling, \textbf{VIN$\rightarrow$E\_GND only}. \\
$C_f$   & \SI{100}{\nano\farad}  & \emph{Optional.} PA1 to \texttt{N\_GND}; buffers the ADC sample-and-hold. Not required at \SI{10}{\kilo\ohm}, but harmless insurance. \\
$R_3$   & \SI{1}{\kilo\ohm}      & Marker series protection (Phase 5). \\
$R_4$   & \SI{10}{\kilo\ohm}     & Marker pulldown to \texttt{E\_GND} (Phase 5). \\
\midrule
\multicolumn{3}{l}{\textit{Reference instrument:} bench multimeter, used for calibration only.} \\
\bottomrule
\end{tabularx}
\end{center}

% ---------------------------------------------------------------------------
\section{Connection table}

\begin{center}
\begin{tabularx}{\textwidth}{llXl}
\toprule
\textbf{Nucleo} & \textbf{Arduino hdr} & \textbf{Connects to} & \textbf{Function} \\
\midrule
\texttt{5V}   & CN6-5 & ESP32 \texttt{VIN} (= \texttt{5V}) & Supply rail \\
\texttt{GND}  & CN6-6 & Bottom of $R_s$ / divider bottom   & System ground, ADC reference \\
\texttt{PA0}  & A0    & Top of $R_s$ = \texttt{E\_GND}     & Shunt sense, \SI{200}{kSps} \\
\texttt{PA1}  & A1    & Divider midpoint                    & Rail sense, $V_\text{rail}/2$ \\
\texttt{PA10} & D2    & ESP32 GPIO27 via \SI{1}{\kilo\ohm}  & Phase marker (optional, 5\,V-tolerant) \\
\midrule
\multicolumn{4}{l}{\textbf{ESP32 USB port: not connected.} See \S\ref{sec:rules}.} \\
\bottomrule
\end{tabularx}
\end{center}

\noindent
\textbf{Why PA10 and not PA4.} PA4 is a DAC pin and is \emph{not} \SI{5}{\volt}-tolerant.
The marker's logic high arrives at up to \SI{3.6}{\volt} (\SI{3.3}{\volt} referenced to a
ground that is itself floating \SI{0.3}{\volt} up), leaving PA4 no headroom. PA10 is an FT pin.

\medskip
\noindent
\textbf{Why GPIO27 on the DUT.} Not a strapping pin, not a flash pin (GPIO6--11), not
input-only (GPIO34/35/36/39), no default peripheral, no boot-time glitch. GPIO13/25/26/33 are
equivalent alternates. \textbf{Never GPIO12} --- it is the MTDI strapping pin; held high at
reset it selects a \SI{1.8}{\volt} flash supply and the board will not boot.

% ---------------------------------------------------------------------------
\section{Electrical design budget}

\subsection{Shunt sizing and LDO headroom}

At the WiFi TX peak, $I_\text{pk} \approx \SI{300}{\milli\ampere}$:
\[
V_{R_s} = I_\text{pk} R_s = \SI{0.3}{\volt},
\qquad
P_{R_s} = I_\text{pk}^2 R_s = \SI{90}{\milli\watt} \;\;(<\;\SI{250}{\milli\watt}\text{ rating}).
\]
The Nucleo's \SI{5}{\volt} pin, when the board is USB-powered, sits near
\SI{4.7}{\volt} (ST-Link protection diode drop). The DUT's AMS1117 therefore sees
\[
V_\text{LDO,in} = \SI{4.7}{\volt} - \SI{0.3}{\volt} = \SI{4.4}{\volt},
\]
which is at --- not above --- the AMS1117's worst-case dropout floor for a
\SI{3.3}{\volt} output.

\medskip
\noindent
\fcolorbox{warn}{white}{\parbox{0.97\textwidth}{
\textbf{Do not increase $R_s$ beyond \SI{1}{\ohm} without addressing headroom.} At
\SI{2}{\ohm} the peak drop is \SI{0.6}{\volt} and the LDO browns out mid-OTA. If the DUT
resets or resets intermittently during WiFi TX, the remedies in order of preference are:
(i) drop to \SI{0.5}{\ohm} (two \SI{1}{\ohm} in parallel), accepting half the sensitivity;
(ii) fit $C_b$ across VIN--E\_GND; (iii) power the Nucleo from an external \SIrange{7}{12}{\volt}
supply on its own \texttt{VIN}, which yields a solid \SI{5.0}{\volt} rail.}}

\subsection{ADC scaling, sampling budget, and why the divider is \SI{1}{\kilo\ohm}}
\label{sec:adc}

Both channels are read by the F446's \SI{12}{bit} SAR ADC against $V_\text{DDA}\approx\SI{3.3}{\volt}$.

\begin{center}
\begin{tabular}{llll}
\toprule
& \textbf{PA0 (shunt)} & \textbf{PA1 (rail)} \\
\midrule
Input range at ADC          & \SIrange{0}{0.5}{\volt} (typ.)  & $\approx\SI{2.35}{\volt}$ \\
Corresponding quantity      & \SIrange{0}{500}{\milli\ampere} & $\approx\SI{4.7}{\volt}$ rail \\
Full-scale before clipping  & \SI{3.3}{\ampere}               & \SI{6.6}{\volt} rail \\
Raw LSB                     & \SI{0.806}{\milli\volt} $\rightarrow$ \SI{0.81}{\milli\ampere} & \SI{1.6}{\milli\volt} rail \\
After $256\times$ decimation & $\approx$ \SI{0.05}{\milli\ampere} & --- \\
Source impedance            & $\approx R_s$, negligible        & \SI{5}{\kilo\ohm} (Th\'evenin, 10k$\|$10k) \\
Sampling time               & 15 cyc.                          & 56 cyc. \\
\bottomrule
\end{tabular}
\end{center}

\noindent
Averaging $N=256$ samples recovers $\log_4 N = 4$ effective bits, giving an effective
\SI{16}{bit} resolution --- roughly \SI{0.05}{\milli\ampere} per count. This is what makes a
bare \SI{12}{bit} ADC with no front-end usable.

\subsubsection*{Divider impedance: \SI{10}{\kilo\ohm}/\SI{10}{\kilo\ohm} is fine}

The ADC input is an RC: the sample-and-hold capacitor ($C_\text{ADC}\approx\SI{4}{\pico\farad}$)
charges through the source impedance plus the internal mux resistance
($R_\text{ADC}\approx\SI{6}{\kilo\ohm}$). Settling to $\tfrac12$\,LSB at \SI{12}{bit} needs
$\approx\ln(2^{13}) \approx 9\tau$. For a \SI{10}{\kilo\ohm}/\SI{10}{\kilo\ohm} divider
(a \SI{5}{\kilo\ohm} Th\'evenin source):
\[
\tau = (\SI{5}{\kilo\ohm} + \SI{6}{\kilo\ohm})\cdot\SI{4}{\pico\farad} \approx \SI{44}{\nano\second},
\qquad
t_\text{settle} \approx 9\tau \approx \SI{0.4}{\micro\second}
\;\ll\;
\underbrace{\frac{56}{\SI{22.5}{\mega\hertz}} = \SI{2.49}{\micro\second}}_{\text{available at 56 cyc.}}
\]
--- roughly \textbf{6$\times$ margin}. The \num{480}-cycle figure sometimes quoted is the spec for
a $\approx\SI{200}{\kilo\ohm}$ source, which a resistor divider is nowhere near.

\medskip
\noindent
\SI{1}{\kilo\ohm}/\SI{1}{\kilo\ohm} works equally well and has more margin still, at the cost of
$\approx\SI{2.5}{\milli\ampere}$ drawn from a rail that is already tight for the DUT's peaks
(\S6.1). That current returns to \texttt{N\_GND} \emph{without passing through the shunt}, so it
does not corrupt the measurement either way. \textbf{Use whichever you have.} The ratio is
\num{2.0} in both cases, so \texttt{KDIV} is unchanged.

\medskip
\noindent
\textbf{Acceptance check, not an argument.} A source impedance too high for its sampling window
biases the rail reading \emph{low}, silently --- and every joule scales with the rail. So verify
it: compare the rail PA1 reports (\texttt{nff power monitor}) against the multimeter. They must
agree. If PA1 reads low, drop to \SI{1}{\kilo\ohm}/\SI{1}{\kilo\ohm} or fit $C_f$.

\subsection{Sampling budget}

The ADC clock is $\text{PCLK2}/4 = \SI{22.5}{\mega\hertz}$, and a \SI{12}{bit} conversion costs
$(t_\text{samp} + 12)$ ADC cycles. One scan of both ranks is therefore
\[
\underbrace{(15+12)}_{\text{PA0}} + \underbrace{(56+12)}_{\text{PA1}} = \num{95}\ \text{cycles}
= \SI{4.22}{\micro\second},
\]
against the \SI{5.00}{\micro\second} that the \SI{200}{kSps} trigger allows --- about
\SI{16}{\percent} headroom. If that budget is ever violated the \texttt{ovr} counter says so;
it is not left to inference. \emph{Measured on the bench:} \num{2040576} samples in
\SI{10}{\second} ($f_s = \SI{199998}{\hertz}$), $\texttt{ovr} = 0$.

\subsection{$V_\text{DDA}$ tolerance}

The Nucleo's \SI{3.3}{\volt} regulator has a $\pm\SI{2}{\percent}$ tolerance, which appears as a
$\pm\SI{2}{\percent}$ gain error on \emph{both} channels.

\medskip
\noindent
\textbf{This is absorbed by the $R_\text{eff}$ calibration (\S\ref{sec:cal}), not corrected
independently.} Calibration solves $R_\text{eff}$ from a measured current, so any gain error
common to the chain --- resistor tolerance, ADC gain, $V_\text{DDA}$ --- lands in that single
constant. They are not separable, and they do not need to be.

\medskip
\noindent
\textit{Not implemented:} the F446 carries a factory VREFINT constant at \texttt{0x1FFF7A2A} from
which the true $V_\text{DDA}$ could be derived independently
($V_\text{DDA} = \SI{3.3}{\volt}\cdot\text{VREFINT\_CAL}/\text{VREFINT\_DATA}$). The firmware does
\emph{not} read it, because calibration already absorbs the error. It would only earn its keep if
one wanted an \emph{uncalibrated} absolute reading, or to track $V_\text{DDA}$ drift with
temperature.

% ---------------------------------------------------------------------------
\section{Firmware integration mathematics}
\label{sec:math}

Let $k$ index samples taken at $\Delta t = \SI{5}{\micro\second}$ ($f_s = \SI{200}{kSps}$), and
let $q_k$, $u_k$ be the \emph{raw ADC counts} on PA0 and PA1. With $L = V_\text{DDA}/4096$ the
volts-per-count, $K$ the divider ratio (\num{2.0} for a matched pair) and $R_\text{eff}$ the
\emph{calibrated} shunt resistance (\S\ref{sec:cal}) --- not the nominal \SI{1}{\ohm}:
\[
I_k = \frac{q_k L}{R_\text{eff}},
\qquad
V_{\text{board},k} = \underbrace{K\,u_k L}_{\text{rail}} - \underbrace{q_k L}_{\text{ground lift}} .
\]
Energy then collapses into sums the firmware can keep in \emph{exact integers}:
\[
E \;=\; \sum_k V_{\text{board},k}\, I_k \,\Delta t
  \;=\; \frac{L^2}{R_\text{eff}}\cdot\frac{1}{f_s}\;
        \Big(K \textstyle\sum_k u_k q_k \;-\; \sum_k q_k^2\Big),
\qquad
Q \;=\; \frac{L}{R_\text{eff}}\cdot\frac{1}{f_s}\sum_k q_k .
\]

\noindent
So the STM32 accumulates only $\sum q$, $\sum q^2$, $\sum u q$, $\sum u$ and $n$ --- as
\texttt{uint64}, which a \SI{10}{\minute} run at \SI{200}{kSps} does not come close to overflowing
--- and the host multiplies by constants. A float accumulator would drift over $10^7$ samples;
these do not.

\medskip
\noindent
\textbf{The host never integrates.} This is deliberate: a host-side integrator that misses serial
frames silently under-reports energy, and the loss is invisible in the output. Here a missed frame
costs nothing --- the next frame carries the correct cumulative total --- and a missed
\emph{conversion} cannot hide either, because it sets \texttt{ovr} (\S\ref{sec:honesty}).

\subsection{Baseline subtraction}

A command's cost is defined as the energy \emph{above} the cost of merely being powered on:
\[
P_\text{idle} = \frac{E(t_1) - E(t_0)}{t_1 - t_0}\bigg|_\text{DUT idle, WiFi connected},
\qquad
\boxed{\;E_\text{marg} = E_\text{total} - P_\text{idle}\cdot T_\text{cmd}\;}
\]
$E_\text{marg}$ is the reported headline figure. Because the LDO loss, USB--UART bridge and
power LED are present in both terms, they cancel.

% ---------------------------------------------------------------------------
\section{Calibration procedure (Phase 0)}
\label{sec:cal}

\noindent\fcolorbox{warn}{white}{\parbox{0.97\textwidth}{%
\textbf{Do not skip, and repeat after \emph{any} rewiring.} Breadboard contact resistance is
\SIrange{10}{100}{\milli\ohm} \emph{per contact} and sits in series with a \SI{1}{\ohm} shunt.
It is not negligible, and it changes when a wire is re-seated.}}

\medskip
\begin{enumerate}[leftmargin=*]
  \item Wire the rig per the schematic. Flash the meter firmware.
  \item Hang a known resistive load on the DUT's \SI{3.3}{\volt} rail --- \SI{100}{\ohm}
        gives $\approx\SI{33}{\milli\ampere}$.
  \item Read the true current $I_\text{DMM}$ with the multimeter \emph{in series} with that load.
  \item Read the current $I_\text{rep}$ the STM32 reports, assuming nominal $R_s = \SI{1}{\ohm}$.
  \item Solve for the effective shunt and persist it:
        \[ R_\text{eff} = R_\text{nom}\cdot\frac{I_\text{rep}}{I_\text{DMM}} \]
        This single constant absorbs the resistor's tolerance, the ADC gain error, \emph{and}
        the breadboard contact resistance.
  \item Repeat for the divider: measure the rail with the DMM, compare against
        $2\,V_\text{PA1}$, and store the ratio correction.
\end{enumerate}

\medskip
\noindent
\textbf{Acceptance criterion.} After calibration, the STM32 and the multimeter must agree to
within a few percent on a purely resistive load. \emph{If they do not, stop} --- nothing built
on top of this rig is worth anything.

% ---------------------------------------------------------------------------
\section{Host interface}

ST-Link Virtual COM Port, \SI{921600}{baud}, 8N1, line-oriented ASCII.

\begin{center}
\begin{tabularx}{\textwidth}{llX}
\toprule
\textbf{Dir} & \textbf{Command} & \textbf{Meaning} \\
\midrule
H$\rightarrow$M & \texttt{PING}                    & Liveness / identification probe. \\
H$\rightarrow$M & \texttt{INFO}                    & Firmware identity, $f_s$, calibration constants. \\
H$\rightarrow$M & \texttt{ZERO}                    & Reset the accumulators and the clock. Acknowledged. \\
H$\rightarrow$M & \texttt{SNAP}                    & Emit one frame \emph{now}. \\
H$\rightarrow$M & \texttt{STREAM <0|1>}            & Free-run frames at \SI{10}{\hertz} (live \texttt{monitor} view only). \\
H$\rightarrow$M & \texttt{WIRECHECK}               & Actively probe whether PA0/PA1 are connected (\S\ref{sec:honesty}). \\
H$\rightarrow$M & \texttt{CAL <micro\_ohm>}        & Set $R_\text{eff}$. \\
H$\rightarrow$M & \texttt{KDIV <milli>}            & Set the divider ratio $\times 1000$ (\texttt{2000} = any matched pair). \\
H$\rightarrow$M & \texttt{VDDA <mv>}               & Set the assumed ADC reference. \\
M$\rightarrow$H & (frame)                          & One JSON object per line. \\
\bottomrule
\end{tabularx}
\end{center}

\noindent
\textbf{Calibration lives on the host} (\texttt{\textasciitilde/.nff/config.json}) and is pushed
down with \texttt{CAL}/\texttt{KDIV}/\texttt{VDDA} on every connect. The meter is stateless across
resets; there is one source of truth.

\medskip
\noindent
\textbf{Frames are emitted on demand, not free-running.} A meter that chatters continuously leaves
the host unable to say exactly when accumulation stopped relative to the profiled command exiting
--- it would over-attribute up to a frame's worth of idle energy to it. \texttt{measure} therefore
drives \texttt{ZERO} $\rightarrow$ run $\rightarrow$ \texttt{SNAP}.

\begin{verbatim}
{"t":10200,"n":2040576,"ovr":0,"sq":141600000,"sqq":9832000000,
 "suq":4.39e11,"su":6.33e9,"qmax":410,"fs":200000,"vdda":3300,
 "r":1000000,"kdiv":2000}
\end{verbatim}

\noindent
The frame carries \emph{raw integer sums}, not derived joules. Two reasons. It keeps the
accumulation exact --- a float accumulator drifts over $10^7$ samples --- and it means a
\emph{re}-calibration can re-derive the energy of a past run without re-measuring it.

% ---------------------------------------------------------------------------
\section{How the rig refuses to lie}
\label{sec:honesty}

Every guard here exists because the rig produced a confident, plausible, \emph{wrong} number
during bring-up. None is hypothetical, and \texttt{ovr} alone catches none of the last two.

\subsection{\texttt{ovr} --- dropped conversions}

A missed conversion means the accumulators are an under-count. If $\texttt{ovr} > 0$, or the
profiled child process exits non-zero, the host reports \texttt{ok: false} and does \textbf{not}
present a joules figure. A plausible-but-wrong energy number is worse than no number, because it
will be trusted.

\subsection{Window cross-check --- a lost \texttt{ZERO}}

A \texttt{ZERO} that never lands leaves the meter accumulating \emph{from boot}, so the frame
describes the board's uptime rather than the command.

\medskip
\noindent
\fcolorbox{warn}{white}{\parbox{0.97\textwidth}{%
\textbf{Observed:} a \SI{30}{\second} run reported \SI{91}{\second} and \SI{26}{\joule} --- with
$\texttt{ovr}=0$, a perfectly sane \SI{58}{\milli\ampere}, and \texttt{ok: true}. The overrun
counter cannot catch this: no samples were lost. The \emph{window} was wrong.}}

\medskip
\noindent
\texttt{ZERO} is now acknowledged, and the host cross-checks (a) the meter's own elapsed clock
against how long it actually waited, and (b) that the samples delivered match the claimed $f_s$
--- if the ADC is not running at $f_s$, every $\Delta t$, and so every joule, is scaled wrong.

\subsection{\texttt{WIRECHECK} --- a rig wired to nothing}

The nastiest failure, and it defeated two designs before this one.

\medskip
\noindent
\fcolorbox{warn}{white}{\parbox{0.97\textwidth}{%
\textbf{An unconnected ADC pin does not read zero --- it holds charge.} With \emph{nothing
attached at all}, the meter reported \SI{590}{\milli\ampere} and a rail of \SI{4.94}{\volt}. Both
are entirely plausible. Any passive range check waves them straight through.}}

\medskip
\noindent
An internal pull-up/pull-down probe fails too: it reported both pins ``connected'' with nothing
attached, because \textbf{the F446's ADC cannot read a pad while the GPIO is in digital-input
mode} (the analog switch is open), and the pulls are \emph{disabled} in analog mode --- so there
is no configuration in which both work.

\medskip
\noindent
What survives is to \textbf{drive the node, release it to analog, and read it immediately}. A real
source restores its own voltage in nanoseconds (the divider is a few-k$\Omega$ source; the shunt
node is a near short to ground). A floating node keeps the charge. Measured: a floating PA0 stayed
\num{1746} counts off, a floating PA1 \num{2425} counts off; a wired node moves under \num{100}.

\medskip
\noindent
For the divider, both tests must pass: it must \emph{read} like a \SI{5}{\volt} rail \emph{and}
spring back. Either alone can be passed by luck --- the bench's floating PA1 sat at a convincing
\SI{4.94}{\volt} --- but it restored to \SI{0.5}{\volt}, and nothing genuinely connected does that.

\medskip
\noindent
\fcolorbox{warn}{white}{\parbox{0.97\textwidth}{%
\textbf{PA0 is only ever driven LOW.} Driving it high would put the GPIO against the
\SI{1}{\ohm} shunt straight to ground --- $\approx\SI{65}{\milli\ampere}$, well past the pin's
\SI{25}{\milli\ampere} rating. PA1 is safe to drive both ways: the divider's series resistance limits
the current to a few \si{\milli\ampere}.}}

% ---------------------------------------------------------------------------
\section{Capability envelope}

\begin{center}
\begin{tabularx}{\textwidth}{clX}
\toprule
& \textbf{Quantity} & \textbf{Remark} \\
\midrule
\yes & Active-mode current, joules per OTA & The target measurand. \\
\yes & Idle / connected-but-quiet baseline & Required for marginal subtraction. \\
\yes & Peak current ($i_\text{max}$)       & Smoothed if $C_b$ is fitted; $Q$ and $E$ remain exact. \\
\midrule
\no  & Deep sleep                          & \SI{10}{\micro\ampere}$\times$\SI{1}{\ohm} = \SI{10}{\micro\volt}, roughly $100\times$ below the breadboard noise floor. Unrecoverable without a real front-end (INA228, PPK2). \\
\no  & \texttt{nff flash}, \texttt{nff monitor} & Require the DUT's USB, which shorts the shunt. \\
\midrule
\multicolumn{3}{l}{\textit{Caveat:} the measurement is whole-board, not ESP32-die. Constant, so it cancels in $E_\text{marg}$.} \\
\bottomrule
\end{tabularx}
\end{center}

% ---------------------------------------------------------------------------
\section{Error budget}

\begin{center}
\begin{tabularx}{\textwidth}{Xllp{58mm}}
\toprule
\textbf{Source} & \textbf{Raw} & \textbf{Residual} & \textbf{Treatment} \\
\midrule
Shunt tolerance ($\pm\SI{5}{\percent}$)      & \SI{5}{\percent}   & $<\SI{0.5}{\percent}$ & Absorbed by $R_\text{eff}$ calibration. \\
Breadboard contact resistance                 & \SIrange{1}{10}{\percent} & $<\SI{1}{\percent}$ & Absorbed by calibration; \textbf{drifts on re-seat}. \\
$V_\text{DDA}$ regulator ($\pm\SI{2}{\percent}$) & \SI{2}{\percent} & $<\SI{0.5}{\percent}$ & Absorbed by $R_\text{eff}$ calibration (\emph{not} VREFINT --- see \S6.4). \\
Divider ratio                                 & \SI{1}{\percent}   & $<\SI{0.5}{\percent}$ & Calibrated against the DMM. \\
Divider source impedance                      & ---                & $\approx 0$        & \SIrange{1}{10}{\kilo\ohm} both settle with $\approx6\times$ margin at \num{56} cycles (\S\ref{sec:adc}). Verify the rail against the DMM. \\
Shunt tempco (self-heating, \SI{90}{\milli\watt}) & --- & $\approx\SI{1}{\percent}$ & Uncorrected. Carbon film $\approx$\SI{300}{ppm\per\celsius}. \\
ADC INL/DNL                                   & ---                & $\approx\SI{0.2}{\percent}$ & Uncorrected, small. \\
Aliasing of WiFi TX spikes                    & ---                & sample-rate bound & Mitigated by $f_s=\SI{200}{kSps}$; $C_b$ trades peak fidelity for stability. \\
Dropped conversions                           & ---                & \textbf{detected} & \texttt{ovr} counter $\rightarrow$ \texttt{ok: false}. \\
Wrong accumulation window                     & ---                & \textbf{detected} & \texttt{ZERO} ack + host/meter clock cross-check. \\
Rig not actually wired                        & ---                & \textbf{detected} & \texttt{WIRECHECK} drive-and-release probe. \\
\midrule
\textbf{Expected overall} & & $\approx\SIrange{2}{3}{\percent}$ & On active-mode energy, post-calibration. \\
\bottomrule
\end{tabularx}
\end{center}

\noindent
The rig is designed for \emph{relative} fidelity: a \SI{2}{\percent} absolute error is
irrelevant to the regression gate, which compares an OTA against a previously measured OTA
on the same, un-rewired rig.

% ---------------------------------------------------------------------------
\section{Operating rules}
\label{sec:rules}

\begin{enumerate}[leftmargin=*,label=\textbf{R\arabic*.}]
  \item \textbf{The ESP32 must not be plugged into the PC's USB.} A USB cable ties the DUT's
        ground to PC ground, which shorts the shunt out. A reading of $\approx\SI{0}{\milli\ampere}$
        during a known-active phase is the signature of exactly this mistake. It is also the
        reason \texttt{nff flash} and \texttt{nff monitor} cannot be profiled here.
  \item \textbf{Only the Nucleo connects to the PC.}
  \item \textbf{$R_s = \SI{1}{\ohm}$, \SI{0.25}{\watt}.} Do not exceed \SI{1}{\ohm} without
        bulk capacitance --- the DUT's LDO will brown out (\S6.1).
  \item \textbf{Recalibrate after any rewiring.} Contact resistance is in the measurement path.
  \item \textbf{All DUT return currents pass through the shunt.} $C_b$ and the marker pulldown
        return to \texttt{E\_GND}, never \texttt{N\_GND}.
\end{enumerate}

% ---------------------------------------------------------------------------
\section{Verification sequence}

Run in order. Each step gates the next.

\begin{enumerate}[leftmargin=*]
  \item \textbf{Resistive-load check.} \SI{100}{\ohm} on the DUT's \SI{3.3}{\volt} rail.
        \texttt{nff power monitor} must agree with the multimeter to within a few percent.
        \emph{If it does not, nothing else counts.}
  \item \textbf{Baseline.} \texttt{nff power measure -{}-for 30}, DUT idle and WiFi-connected.
        Expect tens of \si{\milli\ampere}. A reading near zero means the shunt is shorted ---
        check R1.
  \item \textbf{The real number.} \texttt{nff power measure -{}-during "nff ota ..."}
        $\rightarrow$ joules per OTA. Confirm $\texttt{ovr} = 0$.
  \item \textbf{Regression gate.} Re-run with \texttt{-{}-max-joules} set just below the measured
        value; confirm exit status 1.
  \item \textbf{Agent loop.} Call the \texttt{power\_measure} MCP tool; confirm the agent
        receives a structured verdict.
\end{enumerate}

\medskip
\noindent
\textbf{Reporting.} State the measured joules-per-OTA plainly, including whether \texttt{ovr}
was clean. If the rig cannot be made to agree with the multimeter, say so --- do not report a
number.

% ---------------------------------------------------------------------------
\section{Measurement log}

\begin{center}
\begin{tabularx}{\textwidth}{llllllX}
\toprule
\textbf{Date} & \textbf{$R_\text{eff}$} & \textbf{$P_\text{idle}$} & \textbf{$E_\text{marg}$} & \textbf{$i_\text{max}$} & \textbf{ovr} & \textbf{Notes} \\
\midrule
2026-07-14 & \bad{\SI{1.0}{\ohm} nom.} & \SI{206}{\milli\watt} & \SI{5.37}{\joule} & \SI{314}{\milli\ampere} & 0 & \SI{1}{\mebi\byte} OTA, \texttt{esp32-nffload}. \num{3.06e6} samples. \bad{Uncalibrated} --- resistor measures \SI{1.1}{\ohm}. \\[2mm]
2026-07-14 & --- & --- & --- & --- & 0 & $f_s$ verified: \num{2040576} samples in \SI{10}{\second} $=$ \SI{199998}{\hertz}. \\[2mm]
2026-07-13 & --- & --- & --- & --- & --- & \texttt{WIRECHECK} bring-up: wired to \emph{nothing}, the rig read a plausible \SI{590}{\milli\ampere} and a \SI{4.94}{\volt} rail. \\[2mm]
\bottomrule
\end{tabularx}
\end{center}

\vspace{4mm}
\hrule
\vspace{2mm}
\noindent\small
\textit{Revision A.1 --- rig \textbf{built, debugged and in service} (2026-07-14). It has produced real numbers: see \texttt{nff-power-characteristics.tex}. \bad{Phase 0 calibration still not performed}, so absolute figures carry \SIrange{10}{30}{\percent};
relative figures are good to a few percent. The measurement log below is real.}

\end{document}
